\section{Background and Related Work}
\label{sec:background}

\subsection{Contextual integrity and its benchmarks}

Contextual integrity~\cite{nissenbaum2004,nissenbaum2009} holds that an
information flow is appropriate or not relative to the norms of the context it
occurs in, described by five parameters: data subject, sender, recipient, data
type, and transmission principle. The framing matters for our design because it
makes appropriateness \emph{relational} rather than a property of the data:
telling one's own manager a negotiation walk-away price is entirely
appropriate, and the identical disclosure to the counterparty is not.

A line of benchmarks operationalizes this for language
models---ConfAIde~\cite{confaide2024},
PrivacyLens~\cite{privacylens2024}, CI-Bench~\cite{cibench2024},
CIMemories~\cite{cimemories2025}, and
PrivaCI-Bench~\cite{privacibench2025}. CI-Work~\cite{ciwork2026} identifies
three limitations it sets out to close: prior work isolates a single
information flow where enterprises have parallel entangled flows; prior utility
metrics measure task completion rather than the precision of separating what
must be conveyed from what must be withheld; and prior contexts are short,
where real enterprise retrieval averages roughly 1007 tokens against 55--106 in
earlier benchmarks.

\subsection{CI-Work: the harness we adopt and the gap we close}
\label{sec:background:ciwork}

Because we reuse CI-Work's metrics verbatim, we state its model precisely. A
user $u$ owns a private unstructured store $D$. An agent $A$ receives an
instruction $I$, uses a toolkit $T$, and produces a trajectory
$H=\{(a_t,o_t)\}$. The cumulative retrieved entries $E=\bigcup_t o_t$ are
partitioned into an \emph{essential} set $\Eess$ needed to fulfill $I$ and a
\emph{sensitive} set $\Esens$ whose disclosure to the recipient violates a
privacy norm. Risk is that the final action $\afin$ discloses any member of
$\Esens$ while conveying $\Eess$. A binary judge
$\Fdisc(\afin,e)\in\{0,1\}$ decides disclosure, implemented as
LLM-as-a-Judge~\cite{zheng2023judging} with 83--91\% agreement with human
labels. Three metrics follow:
\begin{align}
\mathrm{LR} &= \tfrac{1}{|\Esens|}\textstyle\sum_{e\in\Esens}\Fdisc(\afin,e)
  &&\text{(Leakage)}\\
\mathrm{VR} &= \max_{e\in\Esens}\Fdisc(\afin,e)
  &&\text{(Violation)}\\
\mathrm{CR} &= \tfrac{1}{|\Eess|}\textstyle\sum_{e\in\Eess}\Fdisc(\afin,e)
  &&\text{(Conveyance)}
\end{align}
with the goal $\mathrm{LR}\downarrow$, $\mathrm{VR}\downarrow$,
$\mathrm{CR}\uparrow$. Cases are stratified over five organizational flow
directions---Downward, Upward, Lateral, Diagonal, and External---at 20\% each
across 125 seeds. Trajectories are simulated in a
ToolEmu-style~\cite{toolemu2024} sandbox in which sensitive atoms are woven
into mixed artifacts so that entanglement is realistic rather than declared.

The results we treat as the problem statement are: violation between 15.8\% and
50.9\% across nine frontier models; a measured privacy--utility trade-off in
which conveyance correlates \emph{positively} with leakage
($r{=}0.40$, $p{=}0.006$) and violation ($r{=}0.39$, $p{=}0.008$); an
inverse-scaling effect where larger variants convey more and leak more; and
prompt-level defenses that leave violation above 20\% while costing 8--12
points of conveyance.

\paragraph{The gap.} Every one of these metrics is defined over $\Esens$ or
$\Eess$, and both sets are defined by \emph{appropriateness of disclosure}. The
benchmark therefore cannot express, let alone score, the question of whether
the agent should have \emph{trusted} an entry. CI-Work's own qualitative
examples contain the unmeasured failure. In its software-renewal case, a
sensitive procurement chat log proposes presenting competitor quotes as \$90
per seat when the figure actually received was \$95; an agent that repeats the
fabricated number to the vendor has not merely leaked, it has let a
low-integrity source ground a high-stakes external commitment. In its
lease-negotiation case, the walk-away threshold is high on \emph{both} axes,
while an after-hours security-incident chat is low-integrity chatter that is
nonetheless sensitive. We develop both cases fully in \Cref{sec:threat}.

\subsection{Multilevel security: Bell--LaPadula and Biba}

Bell--LaPadula~\cite{belllapadula1976} enforces confidentiality with the
simple security property (no read up) and the $\star$-property (no write
down). Biba~\cite{biba1977} is its dual for integrity, with the simple
integrity axiom (no read down) and the $\star$-integrity axiom (no write up).
Both are lattice models in Denning's sense~\cite{denning1976} and together form
the canonical multilevel pair~\cite{nistir7316}. Enforcing both simultaneously
is known to be restrictive, a tension we inherit and manage explicitly through
the threshold and the abstraction path rather than pretend away.

We use Biba rather than Clark--Wilson~\cite{clarkwilson1987}, the other
classical integrity model, for a specific reason: Clark--Wilson's
well-formed-transaction framing presumes a bounded set of certified programs
operating on constrained data items, whereas our objects are natural-language
propositions of heterogeneous provenance arriving at request time. Biba's
lattice comparison over labels is the formulation that survives that setting,
and it composes directly with the confidentiality lattice we already need. We
return to this choice in \Cref{sec:discussion}.

\subsection{Zero trust architecture}

NIST SP 800-207~\cite{nist800207} defines a zero-trust architecture as a Policy
Engine and Policy Administrator producing decisions enforced at a Policy
Enforcement Point, with no implicit trust conferred by network location;
SP 800-207A~\cite{nist800207a} develops the cloud-native form in which sidecar
proxies act as distributed enforcement points. We adopt the PIP/PDP/PEP
decomposition directly, because it names exactly the three jobs our system
does. Notably, SP 800-207 already warns that software agents authenticating to
zero-trust management components may be \emph{coerced into acting on an
attacker's behalf}---which is CI-Work's user-pressure result seen from the
access-control side, a decade earlier and in a different vocabulary.

Tian and Song~\cite{tiansong2021} supply the specific synthesis we build on:
zero trust as dynamic dual-label Bell--LaPadula plus Biba, with continuous
trust scores assigned to users, terminals, channels, files, and applications,
and with \emph{different weights} for confidentiality and integrity
requirements so that decisions balance both per object. \Cref{tab:tiansong}
states what we take and what we change. The three changes in the right column
correspond to the three weaknesses the survey
literature~\cite{he2022zerotrust,sultana2024zerotrust} already identifies in
their method, which is why we treat addressing them as a contribution rather
than as housekeeping.

\begin{table}[t]
\centering
\caption{What we inherit from Tian--Song~\cite{tiansong2021} and what we
change. The right column addresses weaknesses the survey
literature~\cite{he2022zerotrust,sultana2024zerotrust} already identifies.}
\label{tab:tiansong}
\footnotesize
\begin{tabular}{@{}p{0.46\columnwidth}p{0.46\columnwidth}@{}}
\toprule
\textbf{Inherited} & \textbf{Changed or added} \\
\midrule
Dual Bell--LaPadula + Biba labels on every object &
Objects are \emph{retrieved context entries}, not files on disk \\[2pt]
Continuous trust scores per component &
Components are user, agent, tool/channel, entry, recipient \\[2pt]
Weighted confidentiality/integrity decision &
$\wC$, $\wI$ \emph{derived} from breadth and consequence, not hand-set \\[2pt]
``Never trust, always verify'' per request &
Verification at the agent's \emph{write boundary} \\[2pt]
--- &
An explicit initial-trust rule \\[2pt]
--- &
Quantitative evaluation on a benchmark \\
\bottomrule
\end{tabular}
\end{table}

\subsection{Prompt injection is the same failure}
\label{sec:background:injection}

We state this connection early and directly, because a reader who suspects we
have renamed indirect prompt injection~\cite{greshake2023indirect,perez2022ignore}
is entitled to an answer. Our position is that prompt injection \emph{is} an
integrity violation, and Biba is the pre-existing formalism for it: an injected
instruction is a low-integrity input attempting to influence a high-integrity
action. The reframing is useful rather than cosmetic, because it converts an
open detection problem into a labeling problem. Detecting whether a span is
adversarial is hard; labeling a channel as author-controlled is not, and under
the rule in \Cref{sec:method:downgrade} the ceiling on what an author-controlled
entry can cause follows arithmetically from that label.

Three defense lines are the nearest neighbors to ours and we distinguish them
explicitly. AgentDojo~\cite{debenedetti2024agentdojo} evaluates injection
attacks and defenses for agents but scores task utility and attack success
rather than an information-flow property, and has no confidentiality axis.
Wu et al.~\cite{wu2024system} apply information flow control to indirect
injection, which is the closest work on the defense side; our differentiators
are the dual independent axes---their formulation is integrity-only, so it
cannot express the read-but-never-emit case that motivates our write
boundary---and the action-downgrade primitive.
CaMeL~\cite{camel2025} enforces capability-based constraints derived from a
trusted plan, and is the nearest neighbor to our action downgrade; the
distinction is that CaMeL derives capabilities from the control flow of a
planner, whereas we derive the permitted action from the \emph{integrity labels
of the evidence the conclusion rests on}, which is a data-dependent rather than
plan-dependent bound. The dual-LLM pattern~\cite{willison2023dual} anticipates
our L3 privileged-context split informally; we owe it the comparison and make
it in \Cref{sec:method:ladder}.

Our lineage on the systems side is decentralized information flow
control~\cite{myers1997decentralized} and the operating systems that
made it practical~\cite{zeldovich2006histar,krohn2007flume}. The label
propagation problem we face is theirs, with two differences: our objects are
natural-language propositions rather than typed data, so labels must be
\emph{inferred} and carry inference error; and our declassification step is
paraphrase rather than a trusted declassifier, which is why abstraction quality
becomes a measurable property in \Cref{sec:eval} rather than an assumption.
